\section*{BIÊN BẢN PHỎNG VẤN}

\noindent\textbf{VỊ TRÍ:} GIÁM ĐỐC SẢN PHẨM CAO CẤP (GIẢI PHÁP AI) \\
\textbf{ỨNG VIÊN:} Alex Morgan \\
\textbf{NGƯỜI PHỎNG VẤN:} Sarah Jenkins (Phó Chủ tịch Phụ trách Sản phẩm) \\
\textbf{NGÀY \& GIỜ:} 24 tháng 10 năm 2026 \\
\textbf{VÒNG:} Giai đoạn 2 - Kỹ thuật \& Lãnh đạo

\vspace{0.5em}
\noindent\textbf{Bối cảnh Điều hành:} Phỏng vấn vòng hai nhằm đánh giá sự phù hợp chiến lược, nền tảng kỹ thuật trong phát triển sản phẩm AI và năng lực lãnh đạo của ứng viên để dẫn dắt đội ngũ Nền tảng Đám mây Doanh nghiệp.

\subsection*{PHẦN 1: BỐI CẢNH \& TẦM NHÌN}

\noindent\textbf{Sarah Jenkins} \\
Chào mừng Alex. Cảm ơn bạn đã dành thời gian hôm nay. Để bắt đầu, bạn có thể tóm tắt ngắn gọn kinh nghiệm của mình trong việc dẫn dắt các sáng kiến AI doanh nghiệp và điều gì đã thu hút bạn đến với vị trí này tại NexusTech không?

\vspace{0.5em}
\noindent\textbf{Alex Morgan} \\
Cảm ơn Sarah. Rất hân hạnh được ở đây. Trong sáu năm qua tại CloudScale Inc., tôi đã quản lý các sản phẩm B2B SaaS, đặc biệt tập trung vào bộ công cụ phân tích học máy của chúng tôi. Tôi đã dẫn dắt quá trình chuyển đổi từ phân tích truyền thống sang các mô hình AI dự đoán, giúp tăng 35\% ARR. Tôi đã theo dõi sự mở rộng của NexusTech sang các giải pháp AI doanh nghiệp và rất ấn tượng với sự tập trung của bạn vào hạ tầng có khả năng mở rộng. Tôi thấy sự tương thích mạnh mẽ giữa kinh nghiệm tích hợp doanh nghiệp của tôi và lộ trình phát triển của bạn.

\vspace{0.5em}
\noindent\textbf{Sarah Jenkins} \\
Đó là một thành tích rất ấn tượng. Bạn có thể đưa ra một ví dụ cụ thể về cách bạn ưu tiên các tính năng khi nợ kỹ thuật cạnh tranh trực tiếp với các yêu cầu cấp thiết từ khách hàng không?

\vspace{0.5em}
\noindent\textbf{Alex Morgan} \\
Chắc chắn rồi. Năm ngoái, dịch vụ lập chỉ mục cốt lõi của chúng tôi gặp sự cố do nợ kỹ thuật từ hệ thống cũ, làm giảm tốc độ truy vấn gần 40\%. Cùng lúc đó, ba khách hàng doanh nghiệp lớn đang yêu cầu các tính năng báo cáo tùy chỉnh. Tôi đã áp dụng mô hình RICE có trọng số, phân bổ cố định 20\% nguồn lực kỹ thuật mỗi sprint cho việc tái cấu trúc kiến trúc, đồng thời đàm phán mốc thời gian bàn giao theo từng giai đoạn cho các yêu cầu của khách hàng. Trong vòng hai quý, độ trễ truy vấn đã giảm một nửa và chúng tôi đã bàn giao 90\% bộ tính năng được yêu cầu mà không làm kỹ sư bị quá tải.

\subsection*{PHẦN 2: CHIẾN LƯỢC SẢN PHẨM \& GIẢI QUYẾT VẤN ĐỀ KỸ THUẬT}

\noindent\textbf{Sarah Jenkins} \\
Bạn tiếp cận như thế nào khi thiết kế các tính năng dựa trên AI mà đầu ra của mô hình mang tính xác suất thay vì quyết định luận? Bạn quản lý kỳ vọng của người dùng và trải nghiệm người dùng (UX) xung quanh các lỗi như thế nào?

\vspace{0.5em}
\noindent\textbf{Alex Morgan} \\
Thiết kế cho AI đòi hỏi sự dịch chuyển từ kết quả cố định sang khoảng tin cậy. Đầu tiên, chúng tôi thiết kế sự minh bạch trực tiếp vào UI—hiển thị điểm tin cậy hoặc giải thích lý do tại sao một đề xuất lại được đưa ra. Thứ hai, chúng tôi đảm bảo các phương án dự phòng có sự tham gia của con người (human-in-the-loop) diễn ra mượt mà. Nếu ngưỡng tin cậy giảm xuống dưới 80\%, UI sẽ gợi ý xem xét thủ công nhanh thay vì tự động ra quyết định một cách âm thầm.

\vfill
\noindent\small{Đánh giá Tuyển dụng NexusTech • Ứng viên: Alex Morgan \hfill Trang 1 / 2}
\newpage

\subsection*{PHẦN 2: CHIẾN LƯỢC \& GIẢI QUYẾT VẤN ĐỀ (TIẾP THEO)}

\noindent\textbf{Sarah Jenkins} \\
Làm thế nào để bạn xác minh mức độ phù hợp giữa sản phẩm và thị trường (product-market fit) trước khi đầu tư nguồn lực kỹ thuật lớn vào một mô hình AI mới?

\vspace{0.5em}
\noindent\textbf{Alex Morgan} \\
Tôi ủng hộ mạnh mẽ việc thử nghiệm 'Wizard of Oz' hoặc các bản mẫu thử nghiệm độ trung thực thấp (low-fidelity prototypes). Trước khi xây dựng các mô hình tùy chỉnh hoàn chỉnh, chúng tôi thử nghiệm nhu cầu bằng cách sử dụng API của bên thứ ba hoặc các quy trình bán thủ công để quan sát mức độ tương tác của người dùng. Chúng tôi đo lường thành công bằng các chỉ số định tính then chốt như tốc độ hoàn thành quy trình công việc và mức độ hài lòng định tính, thay vì các chỉ số bề nổi ban đầu.

\subsection*{PHẦN 3: LÃNH ĐẠO LIÊN PHÒNG BAN}

\noindent\textbf{Sarah Jenkins} \\
Hãy mô tả một tình huống mà đội ngũ kỹ thuật không đồng ý với chiến lược sản phẩm của bạn. Bạn đã làm thế nào để thống nhất các đội ngũ?

\vspace{0.5em}
\noindent\textbf{Alex Morgan} \\
Các trưởng nhóm kỹ thuật của chúng tôi muốn xây dựng lại toàn bộ hệ thống xử lý dữ liệu (pipeline) bằng một khung làm việc mới, ước tính mất bốn tháng làm việc mà không mang lại giá trị trực tiếp nào cho người dùng. Tôi đã tổ chức một buổi thảo luận chung, nơi chúng tôi đối chiếu các vấn đề khó khăn của nhóm phát triển với dữ liệu khách hàng rời bỏ dịch vụ. Chúng tôi đã thống nhất về một chiến lược chuyển đổi từng bước, giải quyết các điểm nghẽn tồi tệ nhất trước tiên trong khi vẫn tiếp tục phát hành các tính năng có tác động cao trong mỗi chu kỳ phát hành.

\subsection*{ĐÁNH GIÁ PHỎNG VẤN \& ĐIỂM SỐ ỨNG VIÊN}

\begin{table}[h!]
\centering
\begin{tabular}{|l|p{8cm}|c|}
\hline
\textbf{LĨNH VỰC NĂNG LỰC} & \textbf{TÍN HIỆU / BẰNG CHỨNG GHI NHẬN} & \textbf{ĐÁNH GIÁ} \\ \hline
Năng lực Kỹ thuật & Hiểu biết sâu sắc về quy trình làm việc của AI, tích hợp API và các đánh đổi trong hệ thống. & 4.5 / 5.0 \\ \hline
Chiến lược Sản phẩm & Thể hiện tư duy ưu tiên có cấu trúc bằng cách sử dụng các mô hình (RICE, MVP khám phá). & 4.8 / 5.0 \\ \hline
Quản lý Các bên Liên quan & Tiếp cận thực tế đối với các đánh đổi kỹ thuật và giao tiếp với khách hàng. & 4.2 / 5.0 \\ \hline
Lãnh đạo \& Văn hóa & Có tinh thần hợp tác, dựa trên dữ liệu và thể hiện sự đồng cảm rõ ràng với đội ngũ kỹ thuật. & 4.6 / 5.0 \\ \hline
\end{tabular}
\end{table}

\noindent\textbf{Khuyến nghị của Người phỏng vấn: CHUYỂN TIẾP MẠNH MẼ} \\
Alex đã thể hiện chuyên môn sâu sắc về quản lý sản phẩm B2B cùng khả năng nắm bắt xuất sắc các rào cản của sản phẩm AI. Đề xuất chuyển sang giai đoạn phỏng vấn cuối cùng với Phó Chủ tịch Phụ trách Kỹ thuật và CPO.

\vfill
\noindent\small{Đánh giá Tuyển dụng NexusTech • Ứng viên: Alex Morgan \hfill Trang 2 / 2}
